\documentclass[11pt]{article}
% Спільна преамбула для звітів (XeLaTeX)
\usepackage{fontspec}
\usepackage{polyglossia}
\setmainlanguage{ukrainian}
\setotherlanguage{english}
\setmainfont{DejaVuSerif}[Extension=.ttf,UprightFont=*,BoldFont=*-Bold,ItalicFont=*-Italic,BoldItalicFont=*-BoldItalic]
\setsansfont{DejaVuSans}[Extension=.ttf,UprightFont=*,BoldFont=*-Bold,ItalicFont=*-Oblique,BoldItalicFont=*-BoldOblique]
\setmonofont{DejaVuSansMono}[Extension=.ttf,UprightFont=*,BoldFont=*-Bold,ItalicFont=*-Oblique,BoldItalicFont=*-BoldOblique,Scale=0.9]
\usepackage[a4paper,margin=2.2cm]{geometry}
\usepackage{booktabs}
\usepackage{tabularx}
\usepackage{array}
\usepackage{enumitem}
\usepackage{xcolor}
\usepackage{tcolorbox}
\usepackage{graphicx}
\usepackage{setspace}
\usepackage[hidelinks]{hyperref}
\usepackage{titlesec}
\usepackage{fancyhdr}
\usepackage{multirow}
\usepackage{colortbl}
\setlength{\parskip}{4pt}
\setlength{\parindent}{0pt}
\onehalfspacing

\definecolor{accent}{HTML}{1F4E79}
\definecolor{sgreen}{HTML}{E2F0D9}
\definecolor{wred}{HTML}{FBE5D6}
\definecolor{oblue}{HTML}{DEEBF7}
\definecolor{tyellow}{HTML}{FFF2CC}
\definecolor{lightgray}{HTML}{F2F2F2}

\titleformat{\section}{\large\bfseries\color{accent}}{\thesection}{0.6em}{}
\titleformat{\subsection}{\normalsize\bfseries}{\thesubsection}{0.6em}{}

\pagestyle{fancy}
\fancyhf{}
\lhead{\small ІТ, Модуль 2 — командний проєкт}
\rhead{\small Петров Б., Мекшаков А.}
\cfoot{\thepage}
\renewcommand{\headrulewidth}{0.3pt}

\newcolumntype{L}[1]{>{\raggedright\arraybackslash}p{#1}}
\newcolumntype{Y}{>{\raggedright\arraybackslash}X}

% Титульна сторінка
\newcommand{\maketitlepage}[3]{%
\begin{titlepage}
\centering
{\large Київський національний університет імені Тараса Шевченка\\[2pt]
Факультет комп'ютерних наук та кібернетики\par}
\vspace{3cm}
{\large Дисципліна: \textbf{Інформаційні технології}\par}
\vspace{0.5cm}
{\large Модуль 2 — робота в команді\par}
\vspace{2cm}
{\LARGE\bfseries #1\par}
\vspace{0.6cm}
{\large #2\par}
\vfill
\begin{flushright}
\begin{tabular}{l}
\textbf{Виконали:} студенти 4 курсу, група МІ-41\\[4pt]
Петров Богдан\\
Мекшаков Артемій\\[10pt]
\textbf{Спеціальність:} Комп'ютерні науки\\
\end{tabular}
\end{flushright}
\vspace{1.5cm}
{Київ — #3\par}
\end{titlepage}
}

\geometry{margin=2cm}
\setlength{\parskip}{3pt}
\setstretch{1.12}
\newcommand{\say}[1]{\begin{quote}\itshape «#1»\end{quote}}
\newcommand{\who}[1]{\textcolor{accent}{\textbf{[#1]}}}
\begin{document}

\begin{center}
{\LARGE\bfseries Сценарій захисту етапів 2–3}\\[3pt]
{\small $\approx$ 7 хв. Б = Богдан, А = Артемій. Відкрити: обидва PDF, Notion (сторінка + вкладка Kanban).}
\end{center}

\section*{Етап 2 — Ініціація}

\who{Б} \textbf{Стор. 2 — тема.}
\say{StudyFlow — планувальник дедлайнів для студентських груп. Ключова ідея: староста створює дедлайн один раз — його бачить уся група. Стек Next.js + NestJS + PostgreSQL. Тему з етапу 1 залишили; інтеграції з Google винесли за межі семестру — для двох людей це зайвий ризик.}

\who{Б} \textbf{Стор. 3 — Lean Canvas.} Курсором по блоках 1→9.
\say{Проблема: дедлайни розпорошені між Classroom і чатами, староста дублює вручну. Сегменти — студенти і старости, ранні користувачі — наша група. Ціннісна пропозиція: один екран — усі дедлайни групи. Рішення — групові дедлайни, календар, e-mail-нагадування, канбан. Канали — старости й чати груп. Для студентів безкоштовно, потенційно ліцензія для факультету. Витрати — лише час команди, хостинг безкоштовний. Метрики — активні групи і частка вчасних дедлайнів. Перевага — знаємо структуру навчання плюс мережевий ефект у групі.}

\who{Б} \textbf{Стор. 5 — Roadmap.}
\say{Шість версій на рік. MVP 0.1 до 18.11 — ядро з тестами і Swagger; 0.2 до 01.12 — канбан, нагадування, деплой. Далі 1.0 — пілот, 1.5 — інтеграції Google, 2.0 — PWA і Telegram, 2.5 — масштабування на факультет.}

\who{А} \textbf{Стор. 6 → Notion наживо.}
\say{Простір у Notion: головна сторінка з датами, спринтами і правилами роботи — гілки feature, pull request з рев'ю, Definition of Done. База «Беклог» — 45 карток: статус, епік, MoSCoW, квадрант Ейзенхауера, спринт, виконавець, story points.}
Перемкнутись на Notion: головна → база → \textbf{Kanban} → \textbf{Спринти}.
\say{Процес: Backlog → To Do при плануванні спринта → In Progress → Review, рев'ю робить другий → Done після злиття і CI. У To Do зараз — перший спринт. За видом «Спринти» проводимо щотижневе демо.}

\section*{Етап 3 — Планування}

\who{Б} \textbf{Стор. 3 — Project Overview.}
\say{Цілі за SMART: 18.11 — прототип з 20+ тестами, 01.12 — demo і звіт, лютий — пілот із 70 \% дедлайнів групи в системі. Обмеження — двоє людей, нульовий бюджет, жорсткі дедлайни. Ризик не встигнути закриваємо MoSCoW — Could знімається першим; ризик недоступності учасника — обидва знаємо весь стек і робимо рев'ю один одному.}

\who{Б} \textbf{Стор. 4 — Ґант.}
\say{12 місяців, п'ять блоків — це версії з roadmap. Перший — модуль 2 з віхою «захист 01.12», далі пілот, інтеграції, мобільність, масштабування. Ромби — релізи.}

\who{Б} \textbf{Стор. 8 — RACI} (таблиця на стор. 8, не на 5).
\say{18 завдань, на кожне рівно один Accountable. Я — архітектура, БД, backend, нагадування, unit-тести, CI, Swagger. Артемій — вимоги і беклог, UI, frontend, календар і канбан, E2E-тести, звіти, пілот у групі. Рев'ю — обидва. Викладач консультує і приймає етапи — там він єдиний A.}

\who{А} \textbf{Стор. 6 — Ейзенхауер.}
\say{Важливість — внесок у цінність і вимоги курсу, терміновість — чи треба для найближчого демо. Q1 — 16 елементів, спринти 1–2: авторизація, групи, дедлайни, календар. Q2 — тести, документація, канбан, нагадування. Q3 — звіти і слайди, мінімум зусиль. Q4 — після семестру: Telegram, PWA, масштабування.}

\who{А} \textbf{Стор. 9–11 — Беклог.} Прокрутити.
\say{45 елементів у 9 епіках: 33 user story, 9 технічних, 3 організаційні. Кожен має MoSCoW, квадрант, спринт, виконавця і story points за шкалою 1-2-3-5-8. Must — без чого прототип не працює; Won't — свідомо поза семестром. Та сама база, що в Notion.}

\who{А} \textbf{Стор. 12 — Обсяг семестру.}
\say{30 елементів, 76 SP зі 115. Спринти: 22, 24, 22 і 8 SP. Навантаження рівне: Богдан 34, Артемій 32, спільних 10.}

\who{Б} \textbf{Стор. 14 — DoD.}
\say{Усі Must у Done, 20+ unit-тестів у CI, E2E-сценарій реєстрація → група → дедлайн → виконання, Swagger, UML, публічне demo. Дякуємо.}

\section*{Якщо спитають}
\begin{itemize}[leftmargin=1.2em,itemsep=2pt]
  \item \textbf{Чому вдвох?} — Допускається індивідуально; обсяг 76 SP підібрано під двох, ролі по шарах системи.
  \item \textbf{Чому Notion, а не Jira?} — Одна база → кілька видів (Kanban, спринти, Ейзенхауер) плюс документація поряд; для двох людей зручніше.
  \item \textbf{Story point — це що?} — Відносна складність, не години; еталон CRUD предметів = 2, решту порівнювали.
  \item \textbf{Відмінність від Google Calendar?} — Знає групу, предмет і старосту: дедлайн один для всіх, особиста відмітка виконання, прогрес групи.
  \item \textbf{Не встигнете?} — Знімаються Could, потім Should; Must лишається.
  \item \textbf{Як міряти успіх?} — Активні групи на тиждень і частка вчасних дедлайнів; ціль пілоту — 70 \%.
  \item \textbf{Чому модуль 2 в Ґанті один блок?} — Масштаб місячний, бо горизонт рік; тижнева деталізація — у плані спринтів, стор. 12.
\end{itemize}

\textbf{Цифри:} 45 карток · 76/115 SP · спринти 22/24/22/8 · 18 завдань RACI · 6 версій · 18.11 і 01.12.

\end{document}
